\section{Artifact validation and exact status}
\label{sec:validation}

\subsection{Reference implementation}

The file \code{reference/frnf.py} uses only the Python standard library and exact
\code{Fraction} arithmetic.  It implements:

\begin{itemize}
\item least bounded search;
\item prefix maxima and dual-horizon composition;
\item reverse demand closure on a finite DAG;
\item backward viability and least viable paths;
\item majorized frozen descent and the Euclidean example;
\item exact finite-net minimization;
\item stable-observation checking;
\item exact label agreement under a gap; and
\item the guarded reciprocal error transformer.
\end{itemize}

The eleven regression tests in \code{reference/test\_frnf.py} pass.  They include strict
separation between read and justification horizons, a shared-dependency demand graph,
a viability instance with dead ends, rejection of an empty viable path, and exact
rational error checks.  These tests validate the reference algorithms on the included
instances.  They are not proofs of the generic theorems in \(\iPRA\).

\subsection{Lean artifacts}

The Lean directory has three roles.

\begin{description}[style=nextline,leftmargin=2em]
\item[\code{FRNF/Semantic.lean}]
Small host-level records for dual horizons, static transducers, finite candidate
levels, and exact normalizers.
\item[\code{FRNF/Kernel.lean}]
Repository-aligned source for the proof schedule, QF-loop closure, stable observers,
boundary discharge, exact gap collapse, and native export.  It contains no
\code{axiom} declaration or \code{sorry} placeholder.
\item[\code{FRNF/ObjectTargets.lean}]
An explicit roadmap with the bounded-formula syntax and object theorem endpoints for
Booleanization, demand closure, viability, stable clouds, finite-net minima, and
guarded reciprocals.  It intentionally contains visible \code{sorry} placeholders.
\end{description}

Lean was not installed in the artifact runtime, so none of these new files has been
kernel checked.  In particular, absence of placeholders from \code{Kernel.lean} is not
a compilation claim.  The validation folder records every visible placeholder.

\subsection{Paper build}

The PDF is compiled from the included modular LaTeX source.  The build process runs the
Python tests, compiles the bibliography and document, checks that the PDF has extractable
text, and records hashes.  The final PDF is rendered to images and visually inspected
for clipping, missing glyphs, and broken tables.

\subsection{Claim ledger}

\code{metadata/theorem\_ledger.csv} assigns one of the following statuses to every
major theorem: fully argued, reference-tested, upstream object-checked, source-level
Lean combinator, or Lean target.  \code{metadata/obstacle\_matrix.csv} records the
resolution status and remaining input for each difficulty.  These machine-readable
files are intended to survive migration into the repository and prevent claim drift.

\subsection{Current limitations}

The package does not contain:

\begin{itemize}
\item a compiled generic \(\Prf\,\PRAAx\) Booleanization theorem;
\item an object-checked encoded table/demand/viability library;
\item a new effective height theorem for general Siegel or Faltings;
\item an exact Puiseux selector for arbitrary Cauchy coefficient names;
\item stronger analytic zero-free regions or exceptional-zero estimates; or
\item a proof that every theorem in the packet corpus fits \(\FRNF\).
\end{itemize}

What it does provide is a complete mathematical architecture, precise generic proofs,
executable finite algorithms, exact object endpoints, and a prioritized route to a
checked implementation.
